% GDZ PPN236018647,native534middle–535;HKUST535–536.
\NumberedItem{38}
\indent In mentem mihi venit radices aeq[uationis] \(x^n-1\ [=0]\) ex aeq[uationibus] communes radices habentibus elicere ut adeo plerumque tantum aequationes coefficientibus rationalibus gaudentes resolvi oporteat.
\par\noindent\hfill[1796] Sept. 29. Bruns[vigae\textsuperscript{*)}]
\begingroup\renewcommand{\thefootnote}{[*)}\footnotetext{Zwischen dem 17. und dem 29. September 1796 scheint \textsc{Gauss} von Göttingen nach Braunschweig abgereist zu sein. Es wären dann die Tagebuchnotizen Nr. 4 bis Nr. 37 in Göttingen geschrieben; bei der Nr. 38 beginnt auch andere Tinte.]}\endgroup

\medskip\begingroup\small
\indent Wahrscheinlich hat \textsc{Gauss} hier die Zurückführung der Auflösung der Gleichung \(x^n=1\) auf
\[
x^{a^\alpha}=1,\quad x^{{a'}^{\alpha'}}=1,\quad x^{{a''}^{\alpha''}}=1,\quad\ldotp\,\ldotp\,\ldotp
\]
\noindent im Auge, wenn \(n=a^\alpha.{a'}^{\alpha'}.{a''}^{\alpha''}\ \ldotp\,\ldotp\,\ldotp\) ist und \(a,a',a'',\ldotp\,\ldotp\,\ldotp\) lauter verschiedene Primzahlen bedeuten, vergl. \textit{Disquisitiones arithmeticae}, art. 336. (Werke I, S. 413) und art. 366. (Werke I, S. 463); hier ist auch noch weiter von der Zurückführung der Gleichungen die Rede, deren Grad eine Primzahlpotenz ist, auf solche vom Primzahlgrad, was aber die Kenntnis der Lösungen von \(x^\alpha=1,\ x^{\alpha'}=1,\ x^{\alpha''}=1,\ldotp\,\ldotp\,\ldotp\), also die Einführung von Irrationalitäten erfordert. Vergl. auch den ersten Absatz des art. [2.] des nachgelassenen Bruchstücks \textit{Über die Unzerlegbarkeit der Kreisteilungsgleichung}, oben S. 118\begingroup\renewcommand{\thefootnote}{**)}\footnote{Die Angabe von \textsc{Gauss} oben S. 118 Zeile 7 v. u., die Zahlen \(\beta,\beta',\ldotp\,\ldotp\,\ldotp\) seien Divisoren von
\[
a^{\alpha-1}(a-1),\ {a'}^{\alpha'-1}(a'-1),\ \ldotp\,\ldotp\,\ldotp
\]
trifft nicht zu, es kommen vielmehr für \(\beta\) nur die \(a^{\alpha-1}(a-1)\) Zahlen in Frage, die kleiner als \(a^\alpha\) und zu \(a^\alpha\) teilerfremd sind, und entsprechend für \(\beta',\beta'',\ldotp\,\ldotp\,\ldotp\).}\endgroup, und besonders den art. 238. des aus dem Jahre 1797 stammenden 6. Kapitels der \textit{Analysis Residuorum}, Werke II, S. 199.
\par\noindent\hfill\textsc{Loewy.}
\par\endgroup

\NumberedItem{39}
\indent Aequatio tertii gradus est haec:
\[
x^3+xx-nx+\frac{nn-3n-1-mp}{3}=0,
\]
\noindent ubi \(3n+1=p\) et \(m\) numerus resid[uorum] cubic[orum] similes sui excipientes. Unde sequitur si \(n=3k\), fore \(m+1=3l\); si \(n=3k\pm1\), fore \(m=3l\).
\par\vfill\noindent\textsc{X1.}\hfill64
\clearpage\SourceLeaf{535}\ArithmeticHead{506}{L}{TAGEBUCH MIT ERLÄUTERUNGEN.}
\indent Sive
\[
z^3-3pz+pp-8p-9pm=0.
\]
\noindent Hoc [modo]\(m\) penitus determinatum, \(m+1\) semper \(\square+3\square\).
\par\noindent\hfill[1796] Octob. 1. Bruns[vigae]

\medskip\begingroup\small
\indent Es bedeutet hier \(m\) die Anzahl der Lösungen der Kongruenz
\[
1+g^{3t}\equiv g^{3z}\ (\operatorname{mod.}p),
\]
\noindent also dasselbe, was im art. 358. der \textit{Disquisitiones arithmeticae}, Werke I, S. 445 mit \((\mathfrak K,\mathfrak K)=a-1\) bezeichnet wird\footnote{Die dort mit \(n,m\) bezeichneten Größen heißen in unserer Tagebuchnotiz \(p,n\).}. Setzt man ferner, unter Beibehaltung der dortigen Zeichen \(k,a,b,c,\ k=2a-b-c\), so ist
\[
m=\frac{k+n}{3}-1.
\]
\noindent Die erste der obigen Gleichungen nimmt dann die Gestalt an
\[
x^3+x^2-nx-\tfrac19(kp+n)=0
\]
\noindent und geht durch die Substitution \(z=3x+1\) in die zweite Gleichung über. Da zudem nach art. 358. der \textit{Disquisitiones arithmeticae}, Werke I, S. 447, Gleichung II,
\[
m+1=\frac{k+n}{3}=a^2+b^2+c^2-ab-ac-bc
\]
\noindent ist, so folgt
\[
\begin{aligned}
4(m+1)&=(2a-b-c)^2+3(b-c)^2,\\
&=(2b-a-c)^2+3(a-c)^2\\
&=(2c-a-b)^2+3(b-c)^2.
\end{aligned}
\]
\noindent Da nach Werke I, S. 447, Gleichung I, \(a+b+c=n\), also gerade ist, so ist eine der Zahlen \(a,b,c\), etwa \(a\), und die Summe und die Differenz der beiden andern, \(b\pm c\), gerade, also
\[
m+1=\left(a-\frac{b+c}{2}\right)^2+3\left(\frac{b-c}{2}\right)^2.
\]
\noindent Die zweite Gleichung hat \textsc{Gauss} auch in sein Exemplar von \textsc{Lamberts} \textit{Tabellen}, S. 223 eingeschrieben, und zwar in der Form
\[
z^3-3pz-p(9a-(p+1))\ [=0]
\]
\noindent mit der Angabe
\[
4p=3NN+(9a-p-1)^2;
\]
\noindent hier haben \(p,a\) die obige Bedeutung und man hat, wie aus den sieben letzten Zeilen von S. 447, Werke I hervorgeht, \(N=3(b-c)\) zu nehmen. Beispiele findet man in der oben S. 111 abgedruckten Leisteaufzeichnung. Vergl. auch die Tagebuchnotiz Nr. 67 vom 20. Juli 1797, ferner den Artikel 15 des \textsc{Bachmann}schen Aufsatzes, Werke X2, S. 39.

\indent Im Texte muß es statt \textls[120]{excipientes} wohl heißen \textls[120]{excipientium}.
\par\noindent\hfill\textsc{Klein. Bachmann.}
\par\endgroup
